%% ma: L10-B12-0015
%% lop: 10
%% chuong: 5
%% bai: 12
%% dang: 10.12.6
%% loai: TL
%% muc: TH
%% dap_an: "-"
%% nguon: "chuyen-de-cac-so-dac-trung-cua-mau-so-lieu-khong-ghep-nhom-toan-10-chuong-trinh-moi.pdf (GV. Phan Nhật Linh) (P01), Bài 01 Bài tập tự luận 15"
%% kien_thuc: "Lớp 10 Bài 12 (số gần đúng, sai số tuyệt đối, độ chính xác, sai số tương đối, quy tròn số)"
%% trang_thai: chinh-thuc
%% ghi_chu: "Đáp án gốc ý c) là 17660 sai (quy tròn đến hàng chục); theo quy tắc phải quy tròn đến hàng trăm, được 17700 (khớp TN Câu 22 cùng tài liệu). Tự duyệt theo yêu cầu giáo viên 22/10/2026."
\begin{ex}
Viết số quy tròn của mỗi số gần đúng sau với độ chính xác $d$ cho trước:

a) $a=2\,851\,275$ với độ chính xác $d=300$;

b) $a=5{,}2463$ với độ chính xác $d=0{,}001$;

c) $\overline a=17\,658\pm16$.
\loigiai{a) $d=300$: chữ số khác $0$ đầu tiên của $d$ ở hàng trăm, nên quy tròn $a$ đến hàng nghìn: $2\,851\,000$.

b) $d=0{,}001$: hàng phần nghìn, nên quy tròn đến hàng phần trăm: $5{,}25$.

c) $d=16$: hàng chục, nên quy tròn đến hàng trăm: $17\,700$.}
\end{ex}
